\section{Conclusion}
\label{sec:conclusion}

We tested whether a dense residual-stream direction that reads human versus machine authorship also writes AI-like style. In Qwen2.5-7B base and instruct checkpoints, the selected direction achieves 0.9997 held-out AUROC. Adding it during generation, however, produces small, uncertain, and detector-inconsistent changes; nuisance residualization, matched controls, and a reduced base replication do not alter that conclusion. Endpoint quality measures show no broad collapse that would trivially mask the effect.

The key takeaway is that linear readability is not evidence of causal writability. A powered follow-up should use 64--100 held-out prompts per condition, select causal layers on a separate generation-validation set, intervene across adjacent layers, compare addition with projection removal, and obtain blinded human or current-model judgments. Replication in a second model family and directions fit from the target model's own outputs would separate cross-generator authorship structure from write-side alignment. Sparse-feature decomposition is most informative after establishing that a dense causal effect exists.

